\subsection{Application to algebras}

Let A be a ring, B an A-algebra (not necessarily associative or commutative and not necessarily possessing a unit element) and S a multiplicative subset of A. We know that it is possible to define on the \(S^{-1}A\) -module \(S^{-1}B = B \otimes_{A} S^{-1}A\) a canonical \(S^{-1}A\) -algebra structure, said to be obtained by extension of the ring of scalars to \(S^{-1}A\) (Algebra, Chapter III, § 3) and under which the product \((x/s)(y/t)\) is equal to \((xy)/(st)\) . If e is the unit element of B, e/1 is the unit element of \(S^{-1}B\) and if B is associative (resp. commutative), so is \(S^{-1}B\) .

Let A be a ring and M an A-module; we denote by T(M) (resp. ∧(M), S(M)) the tensor algebra (resp. exterior algebra, symmetric algebra) of M (Algebra, Chapter III). It is known that for every commutative A-algebra C there exists a unique isomorphism \(j\) of \(\mathsf{T}(\mathbf{M}) \otimes_{\mathbf{A}} \mathbf{C}\) onto \(\mathsf{T}(\mathbf{M} \otimes_{\mathbf{A}} \mathbf{C})\) (resp. of \(\bigwedge (\mathbf{M}) \oplus_{\mathbf{A}} \mathbf{C}\) onto \(\bigwedge (\mathbf{M} \otimes_{\mathbf{A}} \mathbf{C})\) , of \(\mathsf{S}(\mathbf{M}) \otimes_{\mathbf{A}} \mathbf{C}\) onto \(\mathsf{S}(\mathbf{M} \otimes_{\mathbf{A}} \mathbf{C})\) ) such that

\[
i (x \otimes 1) = x \otimes 1
\]

for all \(x \in \mathbf{M}\) , \(\mathbf{M}\) being canonically identified with a submodule of \(T(\mathbf{M})\) (resp. \(\bigwedge (\mathbf{M}), S(\mathbf{M})\) ) (loc. cit.). Then it is seen in particular that for every multiplicative subset \(\mathbf{R}\) of \(\mathbf{A}\) there are canonical isomorphisms

\[
\begin{array}{c} \mathbf {R} ^ {- 1} \mathsf {T} (\mathbf {M}) \to \mathsf {T} (\mathbf {R} ^ {- 1} \mathbf {M}), \qquad \mathbf {R} ^ {- 1} \bigwedge (\mathbf {M}) \to \bigwedge (\mathbf {R} ^ {- 1} \mathbf {M}), \\ \mathbf {R} ^ {- 1} \mathsf {S} (\mathbf {M}) \to \mathsf {S} (\mathbf {R} ^ {- 1} \mathbf {M}) \end{array}
\]

which reduce to the identity on \(R^{-1}M\) .

